% theme: cell_division_growth_and_asexual_reproduction
\MajorQuestion{細胞分裂と生物の成長／無性生殖}
\SetRowSpaceQanda

\Instruction{細胞分裂に関する次の問いに答えなさい。}
\QQ{1つの細胞が分かれて、2つの細胞になることを\NamedBlank{cell_division}という。}{細胞分裂}
\QQ{生物の体をつくる細胞の\RefBlank{cell_division}を何というか。}{体細胞分裂}
\QQ{植物では、どの部分で\RefBlank{cell_division}がさかんに行われるか。}{根や茎の先端付近}
\QQ{\RefBlank{cell_division}のとき、核の中に見られるひものようなものを\NamedBlank{chromosome}という。}{染色体}
\QQ{\RefBlank{chromosome}は、何とタンパク質からできているか。}{DNA}

\Instruction{細胞分裂と生物の成長に関する次の問いに答えなさい。}
\QQ{\RefBlank{cell_division}の前に、\RefBlank{chromosome}は何倍になるか。}{2倍}
\QQ{生物の特徴を決める遺伝子は、何の一部か。}{DNA}
\QQ{根の先端付近で、細胞分裂がさかんに行われる部分を\NamedBlank{seityou}という。}{成長点}
\QQ{根の先端で、\RefBlank{seityou}を保護している部分を何というか。}{根冠}
\QQ{\NamedBlank{chromosome}に色をつけるために使う液を何というか。}{酢酸オルセイン溶液}

\Instruction{生殖に関する次の問いに答えなさい。}
\QQ{生物が自分と同じ種類の新しい個体をつくることを\NamedBlank{reproduction}という。}{生殖}
\QQ{雄と雌の生殖細胞の核が結びつくことを\NamedBlank{fertilization}という。}{受精}
\QQ{\RefBlank{fertilization}によって新しい個体をつくる\RefBlank{reproduction}を何というか。}{有性生殖}
\QQ{\RefBlank{fertilization}によらず新しい個体をつくる\RefBlank{reproduction}を\NamedBlank{asexual_reproduction}という。}{無性生殖}
\QQ{\RefBlank{asexual_reproduction}でできた個体は、もとの個体の何をそのまま受けつぐか。}{遺伝子}

\Instruction{無性生殖の種類と特徴に関する次の問いに答えなさい。}
\QQ{体が2つに分かれて2つの個体になる\RefBlank{asexual_reproduction}を何というか。}{分裂}
\QQ{体の一部がふくらんで分離し、新しい個体ができる\RefBlank{asexual_reproduction}を\NamedBlank{syutuga}という。}{出芽}
\QQ{植物の体の一部が成長して新しい個体ができる\RefBlank{asexual_reproduction}を\NamedBlank{eiyou}という。}{栄養生殖}
\QQ{\RefBlank{syutuga}によってふえる生物の例を1つ答えなさい。}{ヒドラ}
\QQ{植物の栽培で\RefBlank{eiyou}を利用すると、もとの個体と同じ何をもつ個体をふやせるか。}{形質}
